% Figure: three mode shapes of a three-mass equal-spring chain.
% Each panel shows the chain with mass displacements indicated by
% horizontal arrows.
\begin{figure}[htbp]
\centering
\begin{tikzpicture}[
  springline/.style={decorate, decoration={zigzag, segment length=4pt,
    amplitude=2pt, pre length=4pt, post length=4pt}, draw=enggray, thick},
  wall/.style={draw=engblack, fill=enggray!40},
  mass/.style={draw=engblack, fill=engblue!20, thick,
    minimum width=10pt, minimum height=10pt},
  disp/.style={draw=engred, very thick, -{Stealth}},
  lab/.style={font=\small\bfseries, anchor=south},
  freqlab/.style={font=\small, anchor=west},
]

% Panel 1: Mode 1 -- in-phase, d1=0.35, d2=0.5, d3=0.35
\begin{scope}[shift={(0,0)}]
  \fill[wall] (-0.3,0.4) rectangle (0,1.4);
  \draw (0,0.4) -- (0,1.4);
  \draw[springline] (0,0.9) -- (1.85,0.9);
  \node[mass, anchor=west] (m1) at (1.85,0.9) {$m_{1}$};
  \draw[springline] (m1.east) -- (4.2,0.9);
  \node[mass] (m2) at (4.2,0.9) {$m_{2}$};
  \draw[springline] (m2.east) -- (6.05,0.9);
  \node[mass, anchor=west] (m3) at (6.05,0.9) {$m_{3}$};
  \draw[springline] (m3.east) -- (7.5,0.9);
  \fill[wall] (7.5,0.4) rectangle (7.8,1.4);
  \draw (7.5,0.4) -- (7.5,1.4);
  \draw[disp] (2.1,0.4) -- (2.6,0.4);
  \draw[disp] (4.2,0.4) -- (4.9,0.4);
  \draw[disp] (6.3,0.4) -- (6.8,0.4);
  \node[lab] at (3.7,1.6) {Mode 1 (slow, in phase)};
  \node[freqlab] at (8.0,0.9) {$\omega_{1} = \sqrt{(2 - \sqrt{2})\,k/m}$};
\end{scope}

% Panel 2: Mode 2 -- centre still, d1=0.45, d2=0, d3=-0.45
\begin{scope}[shift={(0,-2.5)}]
  \fill[wall] (-0.3,0.4) rectangle (0,1.4);
  \draw (0,0.4) -- (0,1.4);
  \draw[springline] (0,0.9) -- (1.95,0.9);
  \node[mass, anchor=west] (n1) at (1.95,0.9) {$m_{1}$};
  \draw[springline] (n1.east) -- (3.7,0.9);
  \node[mass] (n2) at (3.7,0.9) {$m_{2}$};
  \draw[springline] (n2.east) -- (5.05,0.9);
  \node[mass, anchor=west] (n3) at (5.05,0.9) {$m_{3}$};
  \draw[springline] (n3.east) -- (7.5,0.9);
  \fill[wall] (7.5,0.4) rectangle (7.8,1.4);
  \draw (7.5,0.4) -- (7.5,1.4);
  \draw[disp] (2.2,0.4) -- (2.85,0.4);
  \draw[disp] (5.85,0.4) -- (5.20,0.4);
  \node[lab] at (3.7,1.6) {Mode 2 (medium, centre still)};
  \node[freqlab] at (8.0,0.9) {$\omega_{2} = \sqrt{2\,k/m}$};
\end{scope}

% Panel 3: Mode 3 -- alternating, d1=0.35, d2=-0.5, d3=0.35
\begin{scope}[shift={(0,-5.0)}]
  \fill[wall] (-0.3,0.4) rectangle (0,1.4);
  \draw (0,0.4) -- (0,1.4);
  \draw[springline] (0,0.9) -- (1.85,0.9);
  \node[mass, anchor=west] (p1) at (1.85,0.9) {$m_{1}$};
  \draw[springline] (p1.east) -- (3.20,0.9);
  \node[mass] (p2) at (3.20,0.9) {$m_{2}$};
  \draw[springline] (p2.east) -- (6.05,0.9);
  \node[mass, anchor=west] (p3) at (6.05,0.9) {$m_{3}$};
  \draw[springline] (p3.east) -- (7.5,0.9);
  \fill[wall] (7.5,0.4) rectangle (7.8,1.4);
  \draw (7.5,0.4) -- (7.5,1.4);
  \draw[disp] (2.1,0.4) -- (2.6,0.4);
  \draw[disp] (3.85,0.4) -- (3.20,0.4);
  \draw[disp] (6.3,0.4) -- (6.8,0.4);
  \node[lab] at (3.7,1.6) {Mode 3 (fast, alternating)};
  \node[freqlab] at (8.0,0.9) {$\omega_{3} = \sqrt{(2 + \sqrt{2})\,k/m}$};
\end{scope}

\end{tikzpicture}
\caption[Three mode shapes of a three-mass linear spring chain]{The
three vibrational modes of three equal masses $m$ coupled by four
equal springs of stiffness $k$ to walls and to each other. Red arrows
show the (small) displacement of each mass at the instant the mode is
maximally extended. Mode 1 (slow) has all three masses moving in
phase, with the centre slightly more than the outers; the restoring
force is small and the frequency low. Mode 2 (medium) has the centre
mass stationary and the outer masses in antiphase; the centre spring
does not stretch and the centre mass sees zero net force. Mode 3
(fast) has all three masses alternating in sign, with maximum spring
extension between each neighbouring pair; the restoring force is
large and the frequency high. The mode shapes are the eigenvectors of
$\mat{K}/m$ and the squared frequencies are its eigenvalues.}
\label{fig:vol02:ch10:vibration-modes}
\end{figure}
